\bookchapter{Capítulo 3: Conecte Seu Celular à “Linha da Vida” — Redes}

\chapterintro{Sem internet, o smartphone fica “meio burro”. A rede é a linha da vida do celular moderno; entendê-la, usá-la bem e protegê-la é uma lição obrigatória para todo usuário de celular.}

\booksection{Conectado sem Confusão — Distinguindo a “Estrada” e o “Carro” do Wi-Fi e dos Dados Móveis}

Muitas pessoas veem os ícones de Wi-Fi e sinal acesos simultaneamente no topo da tela do celular, mas não entendem a diferença essencial entre os dois. O equívoco mais clássico é:\textbf{estar conectado ao Wi-Fi mas não conseguir carregar páginas}— porque o roteador perdeu a internet e o celular não mudou automaticamente para os dados móveis. Para entender isso completamente, precisamos construir um modelo cognitivo de “estrada e carro”.

\begin{tcbinfo}{Metáfora central: Estrada e Carro}
Imagine a internet como uma cidade. A\textbf{“estrada”}é o caminho entre seu celular e a internet, e o\textbf{“carro”}são os dados trafegando na estrada. Wi-Fi e dados móveis são duas estradas diferentes, mas os dados (o carro) são os mesmos.

\end{tcbinfo}

\booksubsection{Diferença Essencial: Canal Local vs Rodovia Remota}

O Wi-Fi é essencialmente uma\textbf{tecnologia de rede local}; é como uma estrada sem fio pavimentada embaixo do seu prédio, responsável por transmitir dados do roteador para o celular. Mas se a banda larga estiver em atraso ou o cabo de fibra óptica rompido, a estrada ainda está lá, mas a “rodovia” externa está cortada — naturalmente, não se vai a lugar nenhum. Já os dados móveis (rede celular) são transmitidos pelas estações base 4G/5G da operadora, como se você entrasse diretamente na via expressa da operadora; enquanto houver sinal, você está diretamente conectado à internet.

\booksubsection{Custo e Cenário: Mensalidade Fixa vs Cobrança por Uso}

O Wi-Fi geralmente é um investimento único (mensalidade da banda larga); enquanto houver energia elétrica, toda a família usa ilimitadamente, sendo ideal para baixar filmes, jogar online e outros cenários de alto tráfego. Já os dados móveis são cobrados por uso; exceder o plano pode resultar em redução de velocidade ou cobrança extra, sendo mais adequado para tarefas leves em cenários móveis como metrô e viagens de negócios.

\begin{tcbpractice}{Experimente: Verifique seu plano de dados}
Abra [Configurações] > [Rede Móvel] > [Gerenciamento de Dados] no seu celular e veja quanto de dados resta neste mês e quais aplicativos mais consomem. Sabendo disso, você fica tranquilo ao sair.

\end{tcbpractice}

\booksubsection{Considerações de Segurança: Wi-Fi Público é um “Campo de Pesca” Oculto}

Em termos de segurança, os dados móveis geralmente são mais seguros que o Wi-Fi público. Os dados móveis são criptografados pela operadora, enquanto o Wi-Fi público pode de fato ser configurado por hackers como um “ponto de acesso falso” para interceptar tráfego não criptografado. Mas é preciso distinguir: com a ampla adoção do HTTPS hoje, o que um atacante pode interceptar são\textbf{quais sites você visitou}(domínios) e\textbf{conteúdo HTTP não criptografado}; já as mensagens do WeChat, senhas de login de aplicativos bancários e outros dados criptografados por HTTPS, o atacante\textbf{não consegue ler diretamente}— a menos que execute um ataque mais complexo de SSL stripping e o usuário ignore os avisos de certificado. Ainda assim, ao fazer operações de pagamento fora de casa, usar diretamente os dados móveis do celular continua sendo a escolha mais segura.

\textbf{Alerta de Segurança}\par

Em locais públicos como aeroportos, shoppings e cafeterias, não se conecte a redes Wi-Fi gratuitas com nomes suspeitos (como “Free\_WiFi\_2024”). Mesmo que conectado, jamais faça login em apps bancários ou insira senhas nessa rede.

\booksubsection{Diagnóstico pelo Celular: Entenda as “Mensagens Secretas” da Barra de Status}

Como saber qual conexão o celular está usando no momento? Em celulares Huawei, por exemplo, o ícone de\textbf{setas bidirecionais}ao lado do sinal da operadora ou do sinal WLAN na barra de status indica que está usando a respectiva via para navegar. Se ambos estiverem acesos simultaneamente, significa que a aceleração de rede está ativa e o celular está usando ambos os canais ao mesmo tempo.

Para usuários de iPhone, é ainda mais intuitivo: se o ícone de Wi-Fi estiver aceso, está usando Wi-Fi; se não estiver, as barras de sinal no canto superior direito indicam que está usando dados móveis. Alguns celulares Android também exibem “5G” ou “4G” na barra de status, informando diretamente o tipo de rede móvel atual.

\booksection{Rede Móvel (4G/5G) e Dados}

A rede móvel é a base para você se conectar a qualquer hora e em qualquer lugar. Do 2G ao 5G, a velocidade aumentou milhares de vezes, mas o princípio é basicamente o mesmo.

\booksubsection{Padrões de rede:}

\begin{itemize}
  \item \textbf{4G LTE}: Atualmente com a cobertura mais ampla, velocidade suficiente (teoricamente até 100Mbps+), consumo de energia relativamente baixo.
  \item \textbf{5G}: Alta velocidade (teoricamente até 10Gbps), mas cobertura ainda em expansão, alto consumo de energia. Adequado para baixar arquivos grandes e assistir vídeos em alta definição.
  \item \textbf{5G SA vs NSA}: SA é rede autônoma (Standalone), NSA é não autônoma (apoiada em estações base 4G). SA oferece melhor experiência, mas requer suporte da operadora e do celular.
\end{itemize}

\booksubsection{Gerenciamento de dados:}

\begin{itemize}
  \item \textbf{O que são dados móveis}: É a quantidade de dados gerada ao navegar na internet pelo celular. Assistir vídeos e navegar por imagens consomem mais dados; texto e voz são relativamente econômicos.
  \item \textbf{Dicas para economizar dados}: Ative o "Modo de economia de dados" nas configurações; atualize aplicativos ou assista vídeos em alta definição apenas no Wi-Fi.
  \item \textbf{Barras de sinal $\neq$ Velocidade}: Sinal forte significa apenas boa conexão; a velocidade também depende de quão ocupada está a estação base e da distância até ela.
\end{itemize}

\textbf{Aviso Importante: Desative o [Roaming de Dados] antes de viajar ao exterior!}\par

Usar dados móveis no exterior pode gerar contas astronômicas! Antes de viajar, desative o roaming de dados nas configurações ou compre antecipadamente um chip de dados local / plano de roaming internacional.

\booksection{Wi-Fi: Conexão, Compartilhamento e Segurança}

O Wi-Fi é o principal meio de conexão em casa e no escritório, rápido e sem consumir dados móveis. Mas o Wi-Fi também está em constante evolução, com grandes diferenças entre os padrões.

\booksubsection{Conectar ao Wi-Fi:}

Vá para [Configurações] > [WLAN], selecione a rede e insira a senha. Muitos celulares agora suportam conexão por leitura de QR code, ainda mais conveniente.

\booksubsection{Compartilhar senha:}

Tanto Android quanto iOS permitem que celulares já conectados compartilhem rapidamente a senha com amigos. O Android gera um QR code, e o iOS pode compartilhar automaticamente por proximidade.

\booksubsection{Evolução dos padrões Wi-Fi:}

Os padrões Wi-Fi são definidos pelo IEEE, com melhorias significativas a cada geração:

\begin{center}
\begin{tabularx}{\linewidth}{|X|X|X|X|X|}
\hline
\textbf{Padrão} & \textbf{Protocolo} & \textbf{Banda} & \textbf{Taxa Máxima} & \textbf{Características} \tabularnewline
\hline
Wi-Fi 4 & 802.11n & 2,4/5GHz & 600Mbps & Padrão antigo, cobertura razoável \tabularnewline
Wi-Fi 5 & 802.11ac & 5GHz & 3,5Gbps & Atualmente o mais difundido, velocidade suficiente \tabularnewline
Wi-Fi 6 & 802.11ax & 2,4/5GHz & 9,6Gbps & Alto desempenho com múltiplos dispositivos, baixa latência \tabularnewline
Wi-Fi 6E & 802.11ax & 2,4/5/6GHz & 9,6Gbps & Adiciona banda de 6GHz, sem congestionamento \tabularnewline
Wi-Fi 7 & 802.11be & 2,4/5/6GHz & 46Gbps & Transmissão ultrarrápida, agregação de múltiplos enlaces \tabularnewline
\hline
\end{tabularx}
\end{center}

\begin{itemize}
  \item \textbf{Tecnologias-chave do Wi-Fi 6}: OFDMA (Acesso Múltiplo por Divisão de Frequência Ortogonal) permite que vários dispositivos transmitam simultaneamente; MU-MIMO (Múltiplos Usuários, Múltiplas Entradas e Múltiplas Saídas) aumenta a eficiência de concorrência entre múltiplos dispositivos; TWT (Tempo de Despertar Programado) torna os dispositivos mais econômicos em energia.
  \item \textbf{Wi-Fi 6E}: Adiciona a banda de 6GHz ao Wi-Fi 6, oferecendo mais canais e reduzindo interferências. Requer suporte tanto do celular quanto do roteador.
  \item \textbf{Wi-Fi 7}: Introduz MLO (Operação de Múltiplos Enlaces), podendo usar simultaneamente várias bandas para transmitir dados, reduzindo drasticamente a latência. Canais ultralargos de 320MHz proporcionam um salto na taxa de transferência. Plataformas de chip flagship como Snapdragon 8 Elite Gen5, Dimensity 9400 já suportam Wi-Fi 7. Modelos como Xiaomi 17, Huawei Mate 80, Samsung S26 podem usar redes Wi-Fi 7 (requer roteador Wi-Fi 7 compatível).
  \item \textbf{2,4GHz vs 5GHz}: 2,4GHz atravessa melhor paredes, mas é mais lento e tem mais interferência; 5GHz é mais rápido, mas atravessa pior paredes. A banda de 6GHz do Wi-Fi 6E/7 é a mais rápida, mas com a pior penetração em paredes, adequada para uso no mesmo cômodo.
\end{itemize}

\booksubsection{Riscos do Wi-Fi público:}

O Wi-Fi público pode ter hackers espionando dados ou configurando pontos de acesso falsos. Regras:

\begin{itemize}
  \item Evite fazer login em contas bancárias ou inserir senhas.
  \item É melhor usar VPN para criptografar a conexão.
  \item Não se conecte automaticamente a Wi-Fi desconhecido.
  \item Use sites HTTPS (ícone de cadeado na barra de endereço).
\end{itemize}

\booksection{Bluetooth: O Campeão Versátil da Comunicação Sem Fio de Curto Alcance}

O Bluetooth é provavelmente uma das tecnologias sem fio que você mais usa no celular — fones de ouvido sem fio, relógios inteligentes, sistemas automotivos, teclados e mouses, todos dependem dele.

\booksubsection{Princípio técnico:}

O Bluetooth opera na banda ISM de 2,4GHz, usando a técnica de espalhamento espectral por salto de frequência (FHSS) para distribuir os dados em 79 canais com troca rápida de transmissão, evitando interferências de forma eficaz. O alcance de comunicação do Bluetooth é dividido em três classes: Classe 1, cerca de 100 metros; Classe 2, cerca de 10 metros (comum em celulares); Classe 3, cerca de 1 metro.

\booksubsection{Evolução das versões:}

\begin{itemize}
  \item \textbf{Bluetooth 4.0}: Introduziu o Bluetooth de Baixa Energia (BLE), permitindo que dispositivos com baterias pequenas, como pulseiras e monitores cardíacos, também se conectassem sem fio, com autonomia passando de dias para meses.
  \item \textbf{Bluetooth 4.2}: Aumentou a velocidade de transmissão e a proteção de privacidade, com suporte a conexão direta IPv6.
  \item \textbf{Bluetooth 5.0}: Velocidade de transmissão dobrada (2Mbps), alcance de transmissão 4 vezes maior (teórico 300 metros), capacidade de broadcast 8 vezes maior. Suporta áudio duplo (conectar dois fones simultaneamente).
  \item \textbf{Bluetooth 5.1}: Introduziu a função de localização direcional (AoA/AoD), com precisão de posicionamento em nível centimétrico.
  \item \textbf{Bluetooth 5.2}: Novo padrão LE Audio, suporta codec LC3 (melhor qualidade de som, menor consumo de energia), áudio multi-fluxo e áudio broadcast.
  \item \textbf{Bluetooth 5.3/5.4}: Otimização adicional de consumo de energia e estabilidade de conexão, broadcast periódico aprimorado.
\end{itemize}

\textbf{Situação da adoção do LE Audio}: Em 2025-2026, o LE Audio entrou em fase de rápida popularização. Os novos fones de ouvido true wireless das principais marcas (como Huawei FreeBuds Pro 5, Xiaomi Buds 6 Pro, Samsung Galaxy Buds3 Pro, etc.) já suportam LE Audio e codec LC3, oferecendo menor latência e melhor qualidade de som. Ao comprar novos fones, recomenda-se priorizar produtos com suporte a LE Audio.

\booksubsection{Cenários de aplicação:}

\begin{itemize}
  \item \textbf{Transmissão de áudio}: Fones de ouvido sem fio, som automotivo. O LE Audio do Bluetooth 5.2 traz um salto qualitativo em qualidade de som e latência.
  \item \textbf{Transmissão de dados}: Transferir arquivos e fotos a curta distância. Embora não seja tão rápido quanto o Wi-Fi Direct, ganha em simplicidade.
  \item \textbf{Pareamento de dispositivos}: Relógios inteligentes, pulseiras, teclados, mouses, controles de jogos.
  \item \textbf{Serviços de localização}: Beacons Bluetooth para navegação interna e rastreamento de itens (como Apple AirTag).
  \item \textbf{Casa inteligente}: Rede Bluetooth Mesh para controlar lâmpadas, fechaduras, etc.
\end{itemize}

\booksection{Controle Remoto Infravermelho: Transforme Seu Celular em um Controle Remoto Universal}

O controle remoto infravermelho é uma das mais antigas tecnologias de controle sem fio, mas ainda é prático no celular. Alguns celulares Android (como certos modelos da Xiaomi e Huawei) possuem emissor infravermelho embutido, podendo controlar TVs, ar-condicionados, decodificadores e outros eletrodomésticos.

\booksubsection{Princípio de funcionamento:}

O controle remoto infravermelho emite sinais de pulso de luz infravermelha em uma frequência específica (geralmente 38kHz) por meio de um diodo emissor de luz infravermelha. O receptor infravermelho do eletrodoméstico decodifica esses pulsos e executa a operação correspondente. Diferentes marcas usam diferentes protocolos de codificação; os mais comuns são NEC, RC5, RC6, SONY, etc. O aplicativo infravermelho no celular contém uma biblioteca de códigos de várias marcas; basta selecionar a marca e o modelo para compatibilidade.

\booksubsection{Limitações:}

\begin{itemize}
  \item \textbf{Direcionalidade}: O infravermelho se propaga em linha reta; é preciso apontar para o dispositivo, não atravessa paredes.
  \item \textbf{Curto alcance}: Distância efetiva geralmente de 5 a 8 metros.
  \item \textbf{Apple não suporta}: O iPhone não possui emissor infravermelho, portanto não pode usar esta função.
\end{itemize}

\booksection{NFC: A Mágica do “Encostar e Pronto”}

NFC (Near Field Communication, Comunicação por Campo de Proximidade) é uma tecnologia de comunicação sem fio de alta frequência de curtíssimo alcance (até 10 centímetros), operando na frequência de 13,56MHz, com taxa de transmissão máxima de 424Kbps.

\booksubsection{Princípio de comunicação e modos:}

O NFC evoluiu da tecnologia RFID e suporta três modos de operação:

\begin{itemize}
  \item \textbf{Modo Leitura/Escrita}: O telefone lê informações de tags NFC, como ler um adesivo NFC para acionar uma ação.
  \item \textbf{Modo Ponto a Ponto}: Dois dispositivos NFC transferem dados entre si, como o Android Beam (descontinuado).
  \item \textbf{Modo de Emulação de Cartão}: O telefone emula um cartão inteligente. Este é o modo mais comum — cartões de transporte, cartões de acesso e cartões bancários usam este modo.
\end{itemize}

\booksubsection{Mecanismos de Segurança:}

\begin{itemize}
  \item \textbf{Distância ultracurta}: A comunicação dentro de 10 cm já é uma garantia de segurança física, tornando a interceptação remota quase impossível.
  \item \textbf{Chip de segurança}: O telefone possui um chip de segurança dedicado (SE) ou usa HCE (Emulação de Cartão via Host). As chaves de pagamento ficam armazenadas em área segura, impossíveis de extrair mesmo com root.
  \item \textbf{Verificação por impressão digital/reconhecimento facial}: O pagamento requer confirmação biométrica, impossibilitando uso fraudulento mesmo com o telefone perdido.
  \item \textbf{Tokenização}: Apple Pay, Huawei Pay e similares não transmitem o número real do cartão, mas usam tokens descartáveis que não podem ser reutilizados mesmo se interceptados.
\end{itemize}

\booksubsection{Aplicações Típicas:}

\begin{itemize}
  \item \textbf{Cartão de transporte/cartão de acesso}: Grave as informações do cartão físico no telefone e passe o telefone.
  \item \textbf{Pagamento móvel}: Apple Pay, Huawei Pay, etc. Vincule seu cartão bancário ao telefone e pague por aproximação.
  \item \textbf{Pareamento de dispositivos}: Aproxime via NFC para parear rapidamente fones Bluetooth e caixas de som.
  \item \textbf{Verificação de identidade}: Leitura de documento de identidade NFC, verificação de carteira de motorista digital.
\end{itemize}

\booksection{VI. Hotspot Pessoal: Transforme o telefone em roteador Wi-Fi}

Quando você tem muitos dados móveis, pode ativar o hotspot para usar no tablet ou computador.

\begin{itemize}
  \item \textbf{Como ativar}: Ajustes > Hotspot Pessoal / Compartilhamento de rede, ative-o.
  \item \textbf{Observações}: Consome muita bateria e dados. Lembre-se de definir uma senha complexa e limitar o número de dispositivos conectados.
  \item \textbf{Compartilhamento USB}: Conecte ao computador via cabo USB, é mais estável e mais rápido que o hotspot sem fio.
\end{itemize}

\begin{tcbpractice}{Mãos à obra: Configure e gerencie o hotspot}
\begin{enumerate}
  \item Abra [Ajustes], encontre "Hotspot Pessoal" ou "Compartilhamento de rede", ative-o e configure o nome e a senha.
  \item Use outro dispositivo (como um computador) para buscar e conectar-se a este Wi-Fi.
  \item No telefone, veja se consegue gerenciar os dispositivos conectados ou remover dispositivos desconhecidos.
\end{enumerate}

\end{tcbpractice}
